\begin{table}[ht]
\centering
\caption{Experiment 1: single-qubit parameter-recovery relative error, mean state fidelity, and positivity-violation rate (fraction of time points with a negative eigenvalue). Values are mean $\pm$ 95\% CI over 5 seeds. True values: $\omega=3.1416$, $\Omega=0.6283$, $\gamma_1=0.3000$, $\gamma_\phi=0.1500$. The classical least-squares fit uses the exact model structure and is the strongest method on this full-data task; the PINN advantage appears under sparse measurements (Table~\ref{tab:exp2}).}
\label{tab:exp1}
\resizebox{\textwidth}{!}{%
\begin{tabular}{lcccccc}
\toprule
Method & $\omega$ err & $\Omega$ err & $\gamma_1$ err & $\gamma_\phi$ err & Mean fidelity & Pos.\ viol. \\
\midrule
CPTP-PINN (ours) & 0.001 $\pm$ 0.001 & 0.026 $\pm$ 0.013 & 0.079 $\pm$ 0.027 & 0.131 $\pm$ 0.050 & 1.0000 $\pm$ 0.0000 & 0.000 \\
Soft-penalty PINN & 0.158 $\pm$ 0.195 & 0.358 $\pm$ 0.287 & 0.266 $\pm$ 0.288 & 1.455 $\pm$ 1.184 & 0.9507 $\pm$ 0.1163 & 0.004 \\
Unconstrained & 0.157 $\pm$ 0.195 & 0.362 $\pm$ 0.284 & 0.276 $\pm$ 0.283 & 1.439 $\pm$ 1.185 & 0.9511 $\pm$ 0.1165 & 0.006 \\
Classical LSQ & 0.001 $\pm$ 0.001 & 0.016 $\pm$ 0.012 & 0.019 $\pm$ 0.008 & 0.020 $\pm$ 0.019 & 1.0000 $\pm$ 0.0000 & 0.000 \\
\bottomrule
\end{tabular}}
\end{table}
